% Generated draft: requires mathematical and source review.
\documentclass[11pt]{article}
\usepackage[margin=25mm]{geometry}
\usepackage{fontspec}
\setmainfont{Latin Modern Roman}
\usepackage{amsmath,amssymb,mathtools,unicode-math}
\setmathfont{Latin Modern Math}
\usepackage{enumitem,needspace,xurl,hyperref}
\setlength{\parindent}{0pt}
\setlength{\parskip}{5pt}
\setlength{\emergencystretch}{3em}
\newcommand{\N}{\mathbb{N}}
\newcommand{\Z}{\mathbb{Z}}
\newcommand{\Q}{\mathbb{Q}}
\newcommand{\R}{\mathbb{R}}
\newcommand{\C}{\mathbb{C}}
\newcommand{\F}{\mathbb{F}}
\newcommand{\A}{\mathbb{A}}
\newcommand{\PP}{\mathbb{P}}
\newcommand{\E}{\mathbb{E}}
\newcommand{\eps}{\varepsilon}
\newcommand{\dd}{\,\mathrm{d}}
\begin{document}
\section*{UnsolvedMath ID: 1146}
\textbf{Source problem code:} catalogue entry
\subsection{Definitions and mathematical statement}
A closed manifold is a compact manifold without boundary. For a connected manifold \(M\) with basepoint \(x\), write \(\pi=\pi_1(M,x)\). It is aspherical if \(\pi_k(M,x)=0\) for every integer \(k\geq2\). A connected space of the homotopy type of a CW complex is a \(K(\pi,1)\) if its fundamental group is isomorphic to \(\pi\) and all its higher homotopy groups vanish. A group \(N\) is perfect if \(N=[N,N]\), where \([N,N]\) is its commutator subgroup.

Under the connected-manifold interpretation used in the supplied literature triage, the displayed source question asks whether, for every closed connected manifold \(M\),
\[
\left[
\forall k\geq2\quad \pi_k(M,x)=0
\quad\text{and}\quad
\forall N\trianglelefteq\pi\quad
\bigl(N=[N,N]\Longrightarrow N=\{1\}\bigr)
\right]
\Longrightarrow M\text{ is a }K(\pi,1).
\]
The original sentence does not explicitly say connected. If disconnected manifolds are allowed, a single fundamental group and a connected \(K(\pi,1)\) conclusion require additional interpretation; that omission is not silently repaired.

The supplied triage dated 2026-08-17 classifies the connected interpretation as true by the definitions, with the perfect-normal-subgroup condition irrelevant to the stated conclusion. It emphasizes that this is a formulation defect, not a resolution of the classical Whitehead conjecture. The latter concerns subcomplexes of aspherical two-dimensional CW complexes and is a different assertion; it is not substituted for the present source statement.

The registry status is open, whereas the supplied triage labels the written implication solved. The earlier background discusses nonpositively curved manifolds and calls the Poincar\'e conjecture a special case, but supplies no precise implication connecting those remarks to this already-aspherical hypothesis. Those remarks are not adopted as mathematical consequences of the displayed statement.
\subsection{Short English statement}
The written question already assumes that the manifold is aspherical, which for a connected manifold is the \(K(\pi_1,1)\) property. The supplied dated triage treats this as a defective formulation rather than the classical Whitehead conjecture. Connectedness is also implicit rather than explicit in the source.
\subsection{Sources}
\begin{itemize}
\item Local UnsolvedMath registry, permanent ID 1146, source code catalogue entry, ``The Whitehead Conjecture,'' including background attribution to J. H. C. Whitehead and literature triage dated 2026-08-17.
\item Allen Hatcher, \emph{Algebraic Topology}, Section 1.B, online textbook: \url{https://pi.math.cornell.edu/~hatcher/AT/AT.pdf}. The supplied triage cites the definitions of asphericity and \(K(G,1)\).
\item A. M. Mikhovich, ``Rational and p-adic analogues of J. H. C. Whitehead's conjecture,'' \emph{Izvestiya: Mathematics} 89 (2025), 274--318: \url{https://doi.org/10.4213/im9597}. The supplied triage cites this for the different classical subcomplex conjecture and its reported open status.
\end{itemize}
\end{document}
